\bookchapter{الفصل 30: عرض الشاشة والعروض التقديمية والاتصال بالسيارة}

\chapterintro{شاشة الهاتف 7 إنش كحد أقصى، لكن تلفازك 55 إنشاً وبروجيكتورك 100 إنش. إلقاء شاشة الهاتف على الشاشة الكبيرة تجربة في بُعد مختلف. هذا الفصل يشرح بث الشاشة والعرض والاتصال بالسيارة، يمتد بقدرات هاتفك لشاشات أكبر.}

\booksection{واحد: ثلاثة معسكرات للبث اللاسلكي — لا تخلط بينها}

البث اللاسلكي يبدو “شاشة الهاتف على التلفاز”، لكن التقنية الأساسية ثلاثة أنواع، غير متبادلة:

\begin{center}
\begin{tabularx}{\linewidth}{|X|X|X|}
\hline
\textbf{التقنية} & \textbf{من يستخدمها} & \textbf{الميزات} \tabularnewline
\hline
AirPlay & iPhone/iPad/Mac & منظومة Apple، التجربة الأكثر سلاسة. يتطلب Apple TV أو تلفاز يدعم AirPlay \tabularnewline
Miracast & هواتف Android & اتصال Wi-Fi Direct مباشر، لا يحتاج راوتر. "البث اللاسلكي" في Android يستخدم هذا افتراضياً \tabularnewline
DLNA & جميع المنصات & يدفع ملفات الوسائط فقط (فيديو/موسيقى/صور)، لا ينسخ الشاشة. مناسب لدفع الفيديو للتلفاز \tabularnewline
\hline
\end{tabularx}
\end{center}

\booksubsection{طرق التشغيل لكل منصة:}

\begin{itemize}
  \item \textbf{iPhone}: مركز التحكم نسخ الشاشة اختر تلفازك/Apple TV. إذا كان التلفاز لا يدعم AirPlay، اشترِ Apple TV أو عصا تلفاز تدعم AirPlay.
  \item \textbf{أندرويد}: اسحب شريط الإشعارات ابحث عن زر "البث اللاسلكي"/"البث الذكي"/"تفاعل متعدد الشاشات" اختر التلفاز. يجب توصيل الهاتف والتلفاز بنفس Wi-Fi.
  \item \textbf{حل عام}: إذا كان التلفاز قديماً لا يدعم أي بث لاسلكي، اشترِ Chromecast أو Xiaomi Box/Huawei Box ووصله في HDMI التلفاز، يمكنه استقبال البث.
\end{itemize}

\booksection{اثنان: البث السلكي — أكثر استقراراً وصفر تأخير}

البث اللاسلكي مريح لكن به تأخير، الألعاب والعرض تجربة سيئة. البث السلكي يحل هذا:

\begin{itemize}
  \item \textbf{iPhone}: يحتاج كابل تحويل Lightning/USB-C إلى HDMI (حوالي 100-200 يوان). بعد التوصيل تُخرج شاشة الهاتف مباشرة لمنفذ HDMI في التلفاز/البروجيكتور.
  \item \textbf{أندرويد}: بعض الهواتف تدعم USB-C إلى HDMI مباشرة (يحتاج الهاتف دعم DP Alt Mode، معظم الرائدة تدعم). اشترِ كابل USB-C to HDMI.
  \item \textbf{سيناريو غرفة الاجتماعات}: إذا كنت تعرض كثيراً من هاتفك، احتفظ بكابل إلى HDMI في حقيبتك. تلفاز/بروجيكتور غرفة الاجتماعات لديه HDMI تقريباً، وصله ويعمل، أكثر موثوقية من البث اللاسلكي.
\end{itemize}

\booksection{ثلاثة: الهاتف يتحول لحاسوب — وضع سطح المكتب}

بعض هواتف Android تدعم “وضع سطح المكتب” — بعد توصيل الشاشة، يُخرج الهاتف واجهة تشبه سطح المكتب، يمكن استخدام الماوس ولوحة المفاتيح:

\begin{itemize}
  \item \textbf{Samsung DeX}: الوظيفة المميزة لسلسلة Samsung الرائدة. بعد توصيل الشاشة تظهر واجهة سطح المكتب، يمكن تعدد النوافذ، التجربة الأقرب لحاسوب حقيقي. مع لوحة مفاتيح وماوس Bluetooth، يمكن السفر دون حاسوب محمول.
  \item \textbf{وضع سطح المكتب في Huawei}: سلسلة Huawei Mate/Pura بعد توصيل الشاشة لها وظائف مماثلة. تدعم تعدد النوافذ وإدارة الملفات وتطبيقات المكتب.
  \item \textbf{Motorola Ready For}: هواتف Moto الرائدة تدعمه، وظائف مماثلة لـ DeX.
  \item \textbf{سيناريوهات الاستخدام}: مهام خفيفة (رد البريد، تعديل المستندات)، عرض PPT، فيديو على شاشة كبيرة. غير مناسب لـ Excel الثقيل أو التصميم الاحترافي.
\end{itemize}

\booksection{أربعة: الاتصال بالسيارة — الطريقة الصحيحة لاستخدام الهاتف أثناء القيادة}

استخدام الهاتف أثناء القيادة مخالف للقانون، لكن الملاحة والموسيقى والمكالمات ضرورة. أنظمة الاتصال بالسيارة تتيح استخدام وظائف الهاتف الأساسية على شاشة السيارة، والهاتف يبقى جانباً:

\begin{itemize}
  \item \textbf{Apple CarPlay}: الخيار الأول لمستخدمي iPhone. عبر USB أو لاسلكي، تظهر واجهة iPhone مبسطة على شاشة السيارة — الملاحة (Amap/Baidu) والموسيقى والهاتف والرسائل والبودكاست. تحكم صوتي Siri بدون يد طوال الوقت.
  \item \textbf{Android Auto}: الحل المقابل لمستخدمي Android. وظائف مشابهة لـ CarPlay، يدعم Google Maps وSpotify وغيرها. في الصين بسبب قيود خدمات Google، بعض الوظائف غير متاحة.
  \item \textbf{Huawei HiCar}: حل Huawei الخاص، معدل دعم السيارات المحلية عالٍ. بالإضافة للملاحة والموسيقى، مكالمات فيديو Huawei ودمج صور الركن وغيرها.
  \item \textbf{Baidu CarLife}: حل عبر المنصات، iOS وAndroid كلاهما يدعمه. وظائف أقل لكن توافق واسع، كثير من السيارات القديمة تدعم CarLife فقط.
\end{itemize}

\begin{tcbwarn}{مستوى متقدم: استخدم GitHub Pages لبناء مدونة شخصية مجانية}
أياً كان حل الاتصال بالسيارة، استخدم التحكم الصوتي قدر الإمكان أثناء القيادة. اضبط الوجهة للملاحة مسبقاً، اختر قائمة الأغاني مسبقاً، لا تتعامل مع شاشة السيارة أثناء القيادة.

\end{tcbwarn}

\booksection{خمسة: بث غرفة الاجتماعات عملياً — بلا ارتباك}

الاجتماعات السيناريو الأكثر تكراراً للبث، والأكثر عرضة للمشاكل. تجارب عملية:

\begin{itemize}
  \item \textbf{احضر مبكراً للاختبار}: لا تنتظر بدء الاجتماع لربط البث. احضر 10 دقائق مبكراً، اختبر طريقة الاتصال. احتمال فشل بث Wi-Fi أعلى مما تظن.
  \item \textbf{أعد حلاً سلكياً كخطة ب}: ضع كابل HDMI في الحقيبة. عند فشل البث اللاسلكي، السلكي منقذ.
  \item \textbf{أغلق معاينة الإشعارات}: أثناء البث ستظهر رسائل WeChat على الشاشة الكبيرة. قبل البث أغلق "معاينة محتوى إشعار شاشة القفل" في الإعدادات، أو فعّل "وضع العرض" (بعض الهواتف تدعم، يخفي الإشعارات تلقائياً عند البث).
  \item \textbf{استخدم "عرض المستند" بدلاً من "تشغيل الهاتف"}: إذا كنت تعرض PPT، احفظه في الهاتف وافتحه بـ WPS بملء الشاشة، بدلاً من بث سطح مكتب الهاتف بأكمله. هكذا أكثر احترافية وسلاسة.
\end{itemize}

\begin{tcbpractice}{جرّب: أكمل بثاً من الهاتف للتلفاز}
\begin{enumerate}
  \item تأكد من اتصال الهاتف والتلفاز بنفس شبكة Wi-Fi.
  \item افتح صورة أو فيديو على الهاتف.
  \item جرّب البث اللاسلكي (iPhone نسخ الشاشة، Android زر البث اللاسلكي).
  \item إذا فشل، تحقق من دعم التلفاز للبروتوكول، أو جرّب دفع DLNA (زر "بث" في الألبوم/تطبيق الفيديو).
  \item بعد نجاح البث، جرّب عمليات أخرى على الهاتف، وانظر إن كانت الشاشة الكبيرة تعرضها متزامنة.
\end{enumerate}

\end{tcbpractice}

\begin{tcbsummary}{خلاصة الفصل}
\begin{itemize}
  \item ثلاثة معسكرات للبث اللاسلكي: AirPlay (Apple) وMiracast (Android) وDLNA (دفع الوسائط عبر المنصات)
  \item البث السلكي أكثر استقراراً وصفر تأخير، غرفة الاجتماعات تحتاج كابل HDMI
  \item Samsung DeX / وضع سطح المكتب في Huawei يحول الهاتف لحاسوب خفيف
  \item الاتصال بالسيارة: CarPlay / Android Auto / HiCar / CarLife أربعة حلول
  \item قواعد بث غرفة الاجتماعات: اختبار مبكر، خطة ب جاهزة، إغلاق معاينة الإشعارات
\end{itemize}

\end{tcbsummary}
